\section{目标化合物及相关结构}
\label{sec:identity}

\subsection{问题与范围}

比较合成条件之前，须先确定研究对象。现有资料包括含Zn目标相的原始报道、无Zn四方Ba--Y硅酸盐的历史研究，以及仅能提供结构参照的结构相关化合物。原始报道将目标相写作$\mathrm{Ba_5Y_{12}Zn[O(SiO_4)]_8}$，并以单晶X射线衍射确定其结构\cite{ababaikeri2024ba5y12zn}；后两类资料虽在组成或晶体学描述上接近目标相，但不能替代含Zn目标相。本节依次核对原始化学式及其出处、历史命名与现代结构的重新审定，并说明现有结构与相区证据所能支持的结论；具体合成变量见第~\ref{sec:condition-matrix}\CJKecglue{}节。

由此可将“获得了何种相”与“采用了何种条件”两个问题分开。表中为每条记录保留原文化学式、本文的排版形式、结构表征结果和待解决的问题。只有化学式与直接结构证据同时指向目标相的研究，才能作为目标相的直接结构依据。同类Ba--Y硅酸盐系列的资料用于识别误指认、非化学计量和竞争相；结构相关化合物仅提示需核查的局部结构单元，不改变目标相的定义。

\subsection{目标相的文献出处}

在本文核查的文献中，$\mathrm{Ba_5Y_{12}Zn[O(SiO_4)]_8}$仅有一篇直接报道\cite{ababaikeri2024ba5y12zn}。该报道采用开放体系高温溶液法生长晶体，并以单晶X射线衍射鉴定其结构\cite{ababaikeri2024ba5y12zn}。这是目前唯一直接确认目标相的研究，其结果不能与相关化合物的条件组合成“共识配方”。该报道支持目标化学式、路线类型和单晶结构归属，但不能证明跨实验室复现或块体相纯窗口，也未完全厘清氧计量表达式、Zn位点及其与无Zn四方硅酸盐系列的关系。

表~\ref{tab:identity-ledger}\CJKecglue{}分列原文写法与本文的排版形式。后者仅统一符号、下标和括号，不将变组成式改写为定组成式，也不依据元素守恒自行修改已发表的化学式。

\begin{table}[bp]
  \centering
  \caption{\textbf{目标化合物及相关Ba--Y硅酸盐的组成与结构依据。}表中记录按文献对象与目标相的关系分类；无Zn研究和结构相关研究均不构成含Zn目标相的独立复现。}
  \label{tab:identity-ledger}
  \begin{tabular}{P{\dimexpr 0.1865\textwidth-2\tabcolsep\relax}P{\dimexpr 0.0813\textwidth-2\tabcolsep\relax}P{\dimexpr 0.2043\textwidth-2\tabcolsep\relax}P{\dimexpr 0.1393\textwidth-2\tabcolsep\relax}P{\dimexpr 0.3885\textwidth-2\tabcolsep\relax}}
    \toprule
    \thead{化学式（原文 → 本文写法）} & \thead{文献} & \thead{主要结构证据} & \thead{与目标相的关系} & \thead{结论与待研究问题}\\
    \midrule
    $\mathrm{Ba_5Y_{12}Zn}$\allowbreak $\mathrm{[O(SiO_4)]_8}$ & \cite{ababaikeri2024ba5y12zn} & 单晶X射线结构鉴定\cite{ababaikeri2024ba5y12zn} & 目标化合物直接文献 & 结论：含Zn目标相及其合成路线类型；待研究：独立复现、块体相纯窗口及与无Zn系列的对应\\
    \addlinespace[3pt]
    $\mathrm{Ba_{5+x}Y_{13}Si_8O_{41}}$（暂定）（本文：保留变量$x$） & \cite{kolitsch2009crystal,wierzbickawieczorek2017high} & 伴生相报道与系统相形成筛查\cite{kolitsch2009crystal,wierzbickawieczorek2017high} & 同类Ba--Y四方硅酸盐 & 结论：历史命名及组成筛查结果；待研究：精确化学计量、结构模型及与目标相的关系\\
    \addlinespace[3pt]
    $\mathrm{Ba_xY_{26}Si_{16}}$\allowbreak $\mathrm{O_{71+x}}$（本文：保留非化学计量式） & \cite{yamane2024synthesis} & 单晶非化学计量调制结构重新审定\cite{yamane2024synthesis} & 同类Ba--Y四方硅酸盐 & 结论：历史式的现代结构解释；待研究：Zn是否可进入该系列\\
    \addlinespace[3pt]
    $\mathrm{Ba_xY_{26}Si_{16}}$\allowbreak $\mathrm{O_{71+x}}$陶瓷 & \cite{yamane2024microstructure} & 块体相组成与显微结构分析\cite{yamane2024microstructure} & 同类Ba--Y四方硅酸盐 & 结论：非化学计量系列的块体表征需求；待研究：第二相、外来SiO$_2$与本征相区的区分\\
    \addlinespace[3pt]
    $\mathrm{Ba_5Y_{13}}$\allowbreak $\mathrm{[SiO_4]_8O_{8.5}}$ & \cite{gulay2024navigation} & EDS、cRED、同步辐射PXRD与中子精修\cite{gulay2024navigation} & 同类Ba--Y四方硅酸盐 & 结论：独立确认的四方Ba--Y结构与相区参照；待研究：与含Zn目标相的位点对应\\
    \addlinespace[3pt]
    $\mathrm{BaY_{16}Si_4O_{33}}$ & \cite{motozawa2022bay16si4o33} & 单晶结构及孤立正硅酸根单元\cite{motozawa2022bay16si4o33} & 结构相关化合物 & 结论：Ba--Y正硅酸盐的结构参照；待研究：缺乏与含Zn目标相同相的依据\\
    \bottomrule
  \end{tabular}
\end{table}

\subsection{四方硅酸盐系列的重新审定}

历史名称的沿用不等于结构的延续。早期助熔剂法生长研究将$\mathrm{Ba_{5+x}Y_{13}Si_8O_{41}}$记为暂定伴生相\cite{kolitsch2009crystal}；后续的系统助熔剂法研究将其纳入多变量相形成筛查，但对象仍为无Zn的Ba--Y四方硅酸盐\cite{wierzbickawieczorek2017high}。现代单晶研究将相关名义组成重新审定为$\mathrm{Ba_xY_{26}Si_{16}O_{71+x}}$非化学计量调制结构\cite{yamane2024synthesis}；相应的陶瓷研究表明，块体样品还须检查第二相和研磨介质引入的SiO$_2$\CJKecglue{}\cite{yamane2024microstructure}。因此，历史名称须依据结构模型和块体相组成重新核对；上述研究均未直接观察到精确的含Zn目标化学式。

独立相区研究结合EDS、cRED、同步辐射PXRD和中子精修，确认了四方$\mathrm{Ba_5Y_{13}[SiO_4]_8O_{8.5}}$\CJKecglue{}\cite{gulay2024navigation}，为无Zn的Ba--Y硅酸盐提供了有别于暂定命名的结构参照。$\mathrm{BaY_{16}Si_4O_{33}}$的单晶结构则给出了含孤立正硅酸根单元的Ba--Y参照结构\cite{motozawa2022bay16si4o33}，可用于检查多面体与硅酸根的连接方式，但不是含Zn目标相的直接证据。本节据此确定目标相的化学式，其余记录用于说明相关结构和竞争相；这些记录所能支持的论断范围见第~\ref{sec:evidence-distance}\CJKecglue{}节。
